% =============================================================================
% viva.tex --- Viva / Project Defence Notes
% Customer Churn Prediction with Explainable AI
% =============================================================================
% Purpose: Maintain what each member has done, how it was done, and why that
% method was chosen --- for proposal, progress reviews, and final viva.
%
% Status tags:
%   [DONE]     Implemented and evidenced in the repository
%   [PLANNED]  Agreed method; update when artifacts exist
%
% Keep aligned with PLAN.md (intent) and PROGRESS.md (evidence).
% Compile: pdflatex viva.tex   (run twice for TOC)
% =============================================================================

\documentclass[11pt,a4paper]{article}

% Minimal preamble for sparse TeX installs. Add geometry/booktabs/hyperref later if available.
\usepackage{graphicx}

% Margins without geometry.sty
\setlength{\oddsidemargin}{0in}
\setlength{\evensidemargin}{0in}
\setlength{\textwidth}{6.5in}
\setlength{\topmargin}{-0.4in}
\setlength{\headheight}{14pt}
\setlength{\headsep}{16pt}
\setlength{\textheight}{9.0in}
\setlength{\parskip}{0.35em}
\setlength{\parindent}{0pt}

\newcommand{\done}{[DONE]}
\newcommand{\planned}{[PLANNED]}
% Incomplete local TeX fonts: map typewriter to roman
\renewcommand{\texttt}[1]{#1}
\makeatletter
\renewcommand{\verbatim@font}{\normalfont\rmfamily}
\makeatother

\title{\textbf{Viva \& Project Defence Notes}\\[0.35em]
\large Customer Churn Prediction with Explainable AI}
\author{4-Member Academic Project Team\\
UCI Online Retail~II Dataset}
\date{Last updated: 25 September 2026\\
\textit{Completion evidence: \texttt{PROGRESS.md}}}

\begin{document}
\maketitle

\begin{abstract}
This file is the team's viva preparation document. For each member it records
\textbf{what} was (or will be) done, \textbf{how} the work is performed, and
\textbf{why} that method was chosen. Completed work is marked \done{} and must
match repository evidence. Planned work is marked \planned{} and must be updated
when artifacts exist. Use this for proposal defence, progress viva, and final
demonstration.
\end{abstract}

\tableofcontents
\newpage

% -----------------------------------------------------------------------------
\section{Project Overview (Common to All Members)}
% -----------------------------------------------------------------------------

\subsection{Problem statement}
E-commerce businesses lose revenue when customers stop purchasing. Early
identification of likely churners enables targeted retention. This project
builds a temporally valid customer-churn classifier on historical retail
transactions, explains predictions with SHAP, and deploys an interactive
Streamlit demonstration.

\subsection{Objective}
\begin{enumerate}
  \item Convert raw invoice line items into a customer-level churn dataset.
  \item Engineer behavioural / RFM features without target leakage.
  \item Train and compare Logistic Regression, Random Forest, and XGBoost
        under a temporal train/validation/test split.
  \item Explain the final model with SHAP and support findings with statistical
        analysis.
  \item Deploy a Streamlit app showing churn probability, risk category, and
        top contributing factors.
\end{enumerate}

\subsection{Dataset}
\begin{tabular}{|p{0.28\textwidth}|p{0.62\textwidth}|}
\hline
\textbf{Field} & \textbf{Value} \\
\hline
Name & Online Retail~II (UCI Machine Learning Repository) \\
\hline
Local path & \texttt{DataSet/online\_retail\_II.xlsx} \\
\hline
Sheets & \texttt{Year 2009-2010}, \texttt{Year 2010-2011} \\
\hline
Raw span & 2009-12-01 to 2011-12-09 \\
\hline
Columns & Invoice, StockCode, Description, Quantity, InvoiceDate, Price, Customer~ID, Country \\
\hline
Citation & Chen, D.\ (2019). \textit{Online Retail II}. UCI ML Repository. \\
\hline
\end{tabular}

\subsection{Development model}
Sequential module-based development with verifiable individual contribution:

\begin{center}
\texttt{M1 (Data)} $\rightarrow$ \texttt{M2 (EDA/Features)} $\rightarrow$
\texttt{M3 (ML)} $\rightarrow$ \texttt{M4 (XAI + Streamlit)}
\end{center}

Upstream artifacts must not be silently redefined. Changes are logged in
\texttt{PROGRESS.md}.

\subsection{Tools}
Python, pandas, Jupyter / VS~Code / Cursor, GitHub, scikit-learn, XGBoost,
SHAP, Streamlit.

\subsection{Current completion snapshot}
\begin{tabular}{|l|l|p{0.48\textwidth}|}
\hline
\textbf{Module} & \textbf{Owner} & \textbf{Status} \\
\hline
M1 Data \& Preprocessing & Member~1 & Implementation \done{}; Git branch/commits pending \\
\hline
M2 EDA \& Features & Member~2 & \planned{} --- not started \\
\hline
M3 Machine Learning & Member~3 & \planned{} --- not started \\
\hline
M4 XAI \& Deployment & Member~4 & \planned{} --- not started \\
\hline
\end{tabular}

\vspace{0.6em}
\noindent\textit{Do not claim M2--M4 as completed in a viva until
\texttt{PROGRESS.md} shows repository evidence.}

\newpage
% -----------------------------------------------------------------------------
\section{Member~1 --- Data Collection \& Preprocessing}
\label{sec:m1}
% -----------------------------------------------------------------------------

\subsection{Role summary}
Own the raw dataset, cleaning rules, temporal periods, churn definition, and
customer-level labelled dataset handed to Member~2.

\subsection{What has been done \done{}}
\begin{enumerate}
  \item Acquired and stored UCI Online Retail~II at
        \texttt{DataSet/online\_retail\_II.xlsx}.
  \item Documented source and citation in \texttt{docs/dataset\_source.md}.
  \item Inspected structure: both sheets concatenated; about 1{,}067{,}371 raw rows.
  \item Implemented reproducible cleaning in \texttt{src/preprocessing.py}
        and \texttt{notebooks/m1\_preprocessing.ipynb}.
  \item Defined observation / prediction periods and churn criteria
        (\texttt{docs/preprocessing\_decisions.md}).
  \item Generated customer-level labels and validation report under
        \texttt{data/processed/}.
\end{enumerate}

\subsubsection*{Key quantitative results}
\begin{tabular}{|l|r|}
\hline
\textbf{Metric} & \textbf{Value} \\
\hline
Raw rows & 1{,}067{,}371 \\
\hline
Missing Customer~ID dropped & 243{,}007 \\
\hline
Cancellations dropped & 18{,}744 \\
\hline
Invalid qty/price dropped & 71 \\
\hline
Exact duplicates dropped & 26{,}124 \\
\hline
Cleaned transactions & 779{,}425 \\
\hline
Eligible customers & 4{,}231 \\
\hline
Churned ($=1$) / Retained ($=0$) & 2{,}217 / 2{,}014 \\
\hline
Churn rate & $\approx$ 52.4\% \\
\hline
Validation & \texttt{all\_passed = true} \\
\hline
\end{tabular}

\subsubsection*{Primary artifacts}
\begin{itemize}
  \item \texttt{data/processed/transactions\_cleaned.parquet}
  \item \texttt{data/processed/customer\_churn\_labels.parquet} (and \texttt{.csv})
  \item \texttt{data/processed/m1\_validation\_report.json}
\end{itemize}

\subsection{How it was done}
\begin{enumerate}
  \item \textbf{Load:} Read both Excel sheets with pandas/\texttt{openpyxl};
        concatenate; keep the raw file untouched.
  \item \textbf{Clean (ordered pipeline):}
        \begin{enumerate}
          \item Drop rows with missing \texttt{Customer ID}.
          \item Drop cancelled invoices (\texttt{Invoice} starts with \texttt{C}).
          \item Drop rows with \texttt{Quantity} $\le 0$ or \texttt{Price} $\le 0$.
          \item Drop exact duplicate rows.
          \item Cast text/ID columns to string; compute
                \texttt{line\_revenue = Quantity $\times$ Price}.
        \end{enumerate}
  \item \textbf{Label:} Keep customers with at least one cleaned purchase in the
        observation window; set \texttt{churn=1} if they have zero purchases in
        the prediction window, else $0$.
  \item \textbf{Validate:} Unique Customer~IDs, non-null labels, label
        consistency, observation-only dates for \texttt{obs\_*} fields, positive
        monetary/invoice counts.
  \item \textbf{Export:} Write parquet/CSV and JSON report via
        \texttt{python -m src.preprocessing}.
\end{enumerate}

\subsection{Why these methods were chosen}
\begin{description}
  \item[Drop missing Customer~ID]
        Churn is customer-level. Guest checkouts cannot be tracked across time;
        keeping them would invent false customers and corrupt labels.
  \item[Drop cancellations (\texttt{C}-invoices)]
        Cancellations are not completed purchases. Including them would inflate
        activity and confuse repurchase detection in the prediction window.
  \item[Drop non-positive quantity/price]
        Invalid or adjustment lines are not reliable sales evidence.
  \item[Drop exact duplicates]
        Identical duplicate rows are recording artifacts, not extra purchases.
  \item[Observation = 2010-01-01 to 2010-12-31]
        A full calendar year of behaviour inside the dataset span, leaving a
        clear future window for labels.
  \item[Prediction = 2011-01-01 to 2011-06-30]
        A six-month forward horizon is long enough to observe repurchase and
        still fits within data ending 2011-12.
  \item[Binary temporal churn]
        Matches ``Will this active customer return next period?'' Features use
        only observation-period information; the future window is used only for
        the label (leakage-aware design).
\end{description}

\subsection{Likely viva questions (Member~1)}
\begin{itemize}
  \item Why not impute Customer~ID instead of dropping?
  \item Why six months for prediction rather than three or twelve?
  \item How do you know \texttt{obs\_*} fields do not leak future data?
  \item What if Member~2 wants a different churn threshold?
\end{itemize}

\subsection{Still outstanding for full faculty DoD}
\begin{itemize}
  \item Git feature branch \texttt{feature/member-1-preprocessing} and commits
        (not yet in repository history).
  \item Screenshots / demo packaging in \texttt{docs/m1\_faculty\_evidence.md}.
\end{itemize}

\newpage
% -----------------------------------------------------------------------------
\section{Member~2 --- EDA \& Feature Engineering}
\label{sec:m2}
% -----------------------------------------------------------------------------

\subsection{Role summary}
Consume M1's customer-level dataset (and cleaned transactions as needed),
perform EDA, engineer RFM and behavioural features, and hand a modelling-ready
feature table to Member~3 --- \textbf{without silently redefining churn labels}.

\subsection{What is planned / done}
\begin{tabular}{|p{0.2\textwidth}|p{0.7\textwidth}|}
\hline
\textbf{Status} & \textbf{Work item} \\
\hline
\planned{} & Load and validate M1 \texttt{customer\_churn\_labels} \\
\hline
\planned{} & Customer / churn distribution analysis \\
\hline
\planned{} & EDA visualisations \\
\hline
\planned{} & RFM features (Recency, Frequency, Monetary) \\
\hline
\planned{} & Behavioural features (AOV, purchase interval, product diversity, cancellation rate, spending/order trends, etc.) \\
\hline
\planned{} & Feature distribution checks, missing values, outlier review \\
\hline
\planned{} & Feature dictionary + \texttt{data/features/} artifact \\
\hline
\end{tabular}

\vspace{0.5em}
\noindent\textit{Update rows to \done{} with paths and figures when M2 finishes.}

\subsection{How it will be done (agreed method)}
\begin{enumerate}
  \item Load \texttt{data/processed/customer\_churn\_labels.parquet}; verify
        row count, churn rate, and schema against
        \texttt{m1\_validation\_report.json}.
  \item Aggregate from \texttt{transactions\_cleaned.parquet} using
        \textbf{observation-period dates only} when building features.
  \item Compute RFM relative to the end of the observation period
        (2010-12-31), not ``today''.
  \item Create behavioural features listed in \texttt{PLAN.md}.
  \item Document distributions, missingness, and outlier handling.
  \item Save engineered dataset under \texttt{data/features/} and document
        definitions in \texttt{docs/}.
\end{enumerate}

\subsection{Why these methods are chosen}
\begin{description}
  \item[RFM features]
        Standard, interpretable retail behavioural summary; strongly related to
        repurchase likelihood and easy to explain in viva.
  \item[Observation-only feature windows]
        Prevents target leakage: using prediction-period purchases as features
        would invalidate metrics.
  \item[Keep M1 churn labels fixed]
        Sequential ownership and faculty verifiability; label changes require a
        Decision Log entry.
  \item[EDA before modelling]
        Surfaces class imbalance and skewed distributions before Member~3 trains
        models.
\end{description}

\subsection{Likely viva questions (Member~2)}
\begin{itemize}
  \item Define Recency, Frequency, and Monetary for this project.
  \item How do you stop leakage from the prediction period?
  \item Which features look most associated with churn (after results exist)?
\end{itemize}

\newpage
% -----------------------------------------------------------------------------
\section{Member~3 --- Machine Learning}
\label{sec:m3}
% -----------------------------------------------------------------------------

\subsection{Role summary}
Build a leakage-aware predictive model: temporal split, imbalance handling,
train three models, evaluate with agreed metrics, select and save the final
model and preprocessing pipeline for Member~4.

\subsection{What is planned / done}
\begin{tabular}{|p{0.2\textwidth}|p{0.7\textwidth}|}
\hline
\textbf{Status} & \textbf{Work item} \\
\hline
\planned{} & Load/validate M2 feature dataset \\
\hline
\planned{} & Temporal train / validation / test split \\
\hline
\planned{} & Class-imbalance analysis and handling \\
\hline
\planned{} & Train Logistic Regression, Random Forest, XGBoost \\
\hline
\planned{} & Evaluate ROC-AUC, PR-AUC, Precision, Recall, F1, Confusion Matrix \\
\hline
\planned{} & Model comparison and final selection with documented rationale \\
\hline
\planned{} & Save model (and pipeline) under \texttt{models/} \\
\hline
\end{tabular}

\subsection{How it will be done (agreed method)}
\begin{enumerate}
  \item Use a \textbf{temporal} split (not a random shuffle that mixes future
        behaviour into training).
  \item Analyse class prevalence; apply an explicit imbalance strategy
        (class weights, resampling, and/or threshold tuning) and record it.
  \item Train Logistic Regression (interpretable baseline), Random Forest
        (non-linear bagging), and XGBoost (strong tabular booster, SHAP-friendly).
  \item Compare on validation metrics; report the frozen test-set result for the
        chosen model.
  \item Persist artifacts for M4 under \texttt{models/}.
\end{enumerate}

\subsection{Why these methods are chosen}
\begin{description}
  \item[Temporal split]
        Random splits can leak future behaviour and overstate churn performance.
  \item[Multiple model families]
        Shows whether gains come from non-linearity/boosting versus a simple
        baseline.
  \item[ROC-AUC and PR-AUC together]
        ROC-AUC alone can look optimistic under imbalance; PR-AUC focuses on the
        churn class.
  \item[Precision, Recall, F1, Confusion Matrix]
        Connect metrics to false alarms versus missed churners.
  \item[Save the actual final model]
        Member~4 must explain and deploy the same model that was selected.
\end{description}

\subsection{Likely viva questions (Member~3)}
\begin{itemize}
  \item Why not only accuracy?
  \item How did you prevent temporal leakage?
  \item Why was the final model selected (cite your experiment table)?
\end{itemize}

\newpage
% -----------------------------------------------------------------------------
\section{Member~4 --- Explainable AI \& Deployment}
\label{sec:m4}
% -----------------------------------------------------------------------------

\subsection{Role summary}
Explain the final M3 model with SHAP, support findings with statistical
analysis, and deliver a Streamlit application for interactive demonstration.

\subsection{What is planned / done}
\begin{tabular}{|p{0.2\textwidth}|p{0.7\textwidth}|}
\hline
\textbf{Status} & \textbf{Work item} \\
\hline
\planned{} & Load final M3 model (and pipeline) \\
\hline
\planned{} & Global SHAP feature importance \\
\hline
\planned{} & Local / per-customer SHAP explanations \\
\hline
\planned{} & Statistical tests / hypothesis analysis \\
\hline
\planned{} & Streamlit UI: customer info, probability, risk, top factors, SHAP \\
\hline
\planned{} & Integration testing and deployment documentation \\
\hline
\end{tabular}

\subsection{How it will be done (agreed method)}
\begin{enumerate}
  \item Load the \textbf{exact} final artifact from \texttt{models/}.
  \item Choose a SHAP explainer compatible with the final model
        (e.g.\ TreeExplainer for tree/boosting models).
  \item Produce global summaries and local explanations for selected customers.
  \item Run documented statistical tests on key feature--churn relationships.
  \item Build Streamlit under \texttt{app/}: select/input customer $\rightarrow$
        preprocess $\rightarrow$ predict $\rightarrow$ show risk + SHAP drivers.
\end{enumerate}

\subsection{Why these methods are chosen}
\begin{description}
  \item[SHAP]
        Unified attribution for global and local explanations; suitable for
        faculty demonstration and business trust.
  \item[Statistical analysis alongside SHAP]
        SHAP explains the model; tests support whether relationships appear in
        the data.
  \item[Streamlit]
        Fast interactive demo for viva walkthroughs.
  \item[Use M3's final model only]
        Explanations and deployment must match reported metrics.
\end{description}

\subsection{Likely viva questions (Member~4)}
\begin{itemize}
  \item Difference between global and local explanation?
  \item What does a high SHAP value mean for a feature?
  \item How does the app avoid using prediction-period inputs?
\end{itemize}

\newpage
% -----------------------------------------------------------------------------
\section{End-to-End Method Narrative (How the Project Is Done)}
% -----------------------------------------------------------------------------

Use this short storyline in proposal / viva introductions:

\begin{enumerate}
  \item \textbf{Data foundation (M1).}
        Start from UCI Online Retail~II invoices. Clean identity, cancellations,
        invalid lines, and duplicates. Define a past observation window and a
        future prediction window. Label churn as non-repurchase in the future
        window among customers active in the past window.
  \item \textbf{Behavioural representation (M2).}
        Summarise each customer's observation-period behaviour with RFM and
        related features so the model sees patterns, not raw invoice rows.
  \item \textbf{Prediction (M3).}
        Train multiple classifiers under a temporal protocol, handle imbalance,
        and select the best evidenced model.
  \item \textbf{Trust \& delivery (M4).}
        Explain why the model predicts churn (SHAP + statistics) and ship an
        interactive Streamlit demo.
\end{enumerate}

\subsubsection*{Pipeline (textual)}
\begin{verbatim}
Online Retail II
  -> Cleaning & filtering
  -> Observation / prediction periods
  -> Churn labels (customer-level)
  -> EDA & feature engineering
  -> Temporal split & imbalance handling
  -> LR / RF / XGBoost + evaluation
  -> SHAP + statistical analysis
  -> Streamlit deployment
  -> Final documentation
\end{verbatim}

\newpage
% -----------------------------------------------------------------------------
\section{Contribution Matrix (Quick Viva Reference)}
% -----------------------------------------------------------------------------

\begin{tabular}{|p{0.14\textwidth}|p{0.2\textwidth}|p{0.3\textwidth}|p{0.24\textwidth}|}
\hline
\textbf{Member} & \textbf{Owns} & \textbf{Key method} & \textbf{Why} \\
\hline
Member~1 & Labels \& cleaned data & Temporal churn windows + cleaning rules & Leakage-aware, valid customer labels \\
\hline
Member~2 & Features \& EDA & RFM + behavioural feats from observation only & Interpretable, no target leakage \\
\hline
Member~3 & Models \& metrics & Temporal split; LR/RF/XGB; multi-metric eval & Fair comparison, realistic performance \\
\hline
Member~4 & XAI \& app & SHAP + stats + Streamlit & Explainability and demo \\
\hline
\end{tabular}

\vspace{1em}
\noindent\textbf{Evidence habit:} When answering ``What did you do?'', cite a
file path, a metric, or a commit --- not only a verbal claim.

\newpage
% -----------------------------------------------------------------------------
\section{How to Maintain This File}
% -----------------------------------------------------------------------------

\begin{enumerate}
  \item After a module finishes, change \planned{} to \done{} for that member.
  \item Paste key numbers, file paths, and (when available) commit hashes.
  \item Sync claims with \texttt{PROGRESS.md}; never invent completed work.
  \item Add new Decision Log items into the member's ``Why'' subsection.
  \item Recompile: \texttt{pdflatex viva.tex} (twice if the TOC changes).
\end{enumerate}

\subsection{Update log}
\begin{tabular}{|l|l|p{0.55\textwidth}|}
\hline
\textbf{Date} & \textbf{Who} & \textbf{Change} \\
\hline
2026-09-25 & Team & Initial \texttt{viva.tex}: M1 what/how/why from repository evidence; M2--M4 planned methods from \texttt{PLAN.md}. \\
\hline
\end{tabular}

\end{document}
